\documentclass[a4paper, 11pt]{article}

\usepackage[T1,T2A]{fontenc}
\usepackage[utf8]{inputenc}
\usepackage[english, russian]{babel}

\usepackage{comment}
\usepackage{fullpage}
\usepackage{amssymb}
\usepackage{amsmath}
\usepackage{amsthm}
\usepackage{enumitem}

% %\DeclareMathOperator{\ord}{ord}
\DeclareMathOperator{\chr}{char}
\DeclareMathOperator{\Mod}{mod}
\DeclareMathOperator{\Gal}{Gal}
\DeclareMathOperator{\N}{N}
\DeclareMathOperator{\Tr}{Tr}
\DeclareMathOperator{\SL}{SL}

\begin{document}

\begin{center} \Large \bf Листок 5\\ Тема 5 (2.1). Конечные поля \end{center}

\noindent\textbf{Упражнения и задачи}

\begin{enumerate}[topsep=0pt]
    \item Завершите доказательство леммы: пусть $k$ --- поле характеристики $p$, тогда $\forall \alpha, \beta \in k$ $\forall d \in \mathbb{Z}_+$ $(\alpha+\beta)^{p^d}=\alpha^{p^d}+\beta^{p^d}$.
    \item Докажите, что многочлен $x^2+x+1$ неприводим над $\mathbb{F}_2$ (т.е. $\mathbb{F}_2[x]/(x^2+x+1)$ является полем). Выпишите таблицы операций сложения и умножения элементов в этом поле, а также обратные элементы.
    \item Найдите все неприводимые многочлены степени $4$ над полем $\mathbb{F}_2$.
    \item Пусть $p,q$ --- различные простые. Чему равно число неприводимых многочленов степени $q$ в $\mathbb{F}_p[x]$? %[IR, p87 ex7.17]
    \item Докажите следующие оценки для $\nu(n)$ --- числа неприводимых многочленов степени $n\geqslant 2$ над $\mathbb{F}_p$:
    $$
    \tfrac{1}{n} p^n - \tfrac{p}{n(p-1)}(p^{\frac{n}{2}}-1) \leqslant \nu(n) \leqslant \tfrac{1}{n}(p^n-p).
    $$
    \item Пусть $\sigma_j(f) = \sum_{g|f}' (Ng)^j$, где $Nf=q^{\deg f}$, а суммирование берется по всем унитарным многочленам $g(x)$. Докажите, что
    \begin{itemize}[noitemsep,topsep=0pt]
        \item $\sum\limits_f \dfrac{\sigma_0(f)}{(Nf)^s} = \dfrac{1}{(1-q^{1-s})^2}$;
        \item $\sum\limits_f \dfrac{\sigma_1(f)}{(Nf)^s} = \dfrac{1}{(1-q^{1-s})(1-q^{2-s})}$.
    \end{itemize} %[IR, p87 ex7.20]
    \item Пусть $\alpha \in \mathbb{F}_q^*$ ($q=p^s$). Докажите, что $x^n=\alpha$ разрешимо $\Leftrightarrow$ $\alpha^{(q-1)/d}=1$, где $d=(n,q-1)$, причем если разрешимо, то $d$ решений. %[IR, p80 prop 7.1.2]
    \item Как выглядит подгруппа всех квадратов в $\mathbb{F}_{2^n}$? %[IR, p86 ex7.8]
    \item Пусть $n|q-1$, докажите, что $G=\{\alpha\in\mathbb{F}_q^*: x^n=\alpha\ \text{--- разрешимо}\}$ --- подгруппа в $\mathbb{F}_q^*$, $|G|=\frac{q-1}{n}$. %[IR, p86 ex7.4]
    \item Пусть $q=p^n, p>2$. Определим аналог символа Лежандра для $\mathbb{F}_q$: $\left(\frac{\alpha}{q}\right)=1$, если $\alpha$ --- квадрат в $\mathbb{F}_q$; $\left(\frac{\alpha}{q}\right)=-1$, если $\alpha$ не является квадратом в $\mathbb{F}_q$; $\left(\frac{0}{q}\right)=0$. Докажите следующие свойства этого символа: 
    \begin{itemize} [topsep=0pt]
        \item $\left(\frac{\alpha\beta}{q}\right) = \left(\frac{\alpha}{q}\right)\left(\frac{\beta}{q}\right)$, $\alpha, \beta \in \mathbb{F}_q$; 
        \item $\sum_{\alpha\in\mathbb{F}_q} \left(\frac{\alpha}{q}\right) = 0$; 
        \item $\left(\frac{\alpha}{q}\right)=\left(\frac{N_{\mathbb{F}_q/\mathbb{F}_p}(\alpha)}{p}\right)$ --- обычный символ Лежандра $\bmod\, p$.
    \end{itemize} %[Степ Задача I.2.9]
\end{enumerate}

% \noindent\textbf{SageMath}
% \begin{itemize}[topsep=0pt]
%     \item Исследуйте основные функции SageMath связанные с заданием и свойствами конечных полей
%     \begin{itemize}[noitemsep,topsep=0pt]
%         \item Определение конечного поля: \texttt{FiniteField(), GF()};
%         \item Неприводимый многочлен задающий конечное поле: \texttt{polynomial()}, опция \texttt{modulus} в \texttt{FiniteField()} для явного задания неприводимого многочлена модели конечного поля;
%         \item Решение уравнения $x^n=\alpha$: \texttt{nth\_root()};
%         \item Поле разложения: \texttt{splitting\_field()};
%         \item Расширение полей: \texttt{extension()}.
%      \end{itemize}
% \end{itemize}

\noindent\textbf{Темы для самостоятельного изучения}
\begin{itemize}[topsep=0pt]
    \item ---
\end{itemize}


\end{document}
